% Chapter: 05_discussion
% Included from lectures/main.tex

{
\setbeamertemplate{section page}{
  \begin{centering}
    \usebeamerfont{section title}\insertsection\par
  \end{centering}
  \note{
  \begin{itemize}
  \item You now have a rough overview of what happened the last years in MechInterp
  \item Where do we go from here?
  \item Open question: what kind of progress would you want to see in MechInterp? What is realistic?
  \end{itemize}
  }
}
\section{Discussion}
}

\begin{frame}{Open Problems in MechInterp}
Decomposition: problems with Sparse Dictionary Learning
\begin{itemize}
\item High reconstruction errors
\item Expensive for large models
\item Sparsity is a flawed proxy for interpretability
\item Dataset dependent
\item No clarity on the concept of a ``feature''
\end{itemize}
\source{sharkey2025}

\note{
\begin{itemize}
\item Neurons, attention heads, and layers are polysemantic and don't decompose cleanly
\item SDL (SAEs etc.) has high reconstruction errors, is expensive for large models, assumes linearity in nonlinear models, and uses sparsity as a flawed proxy for interpretability
\item SDL leaves feature geometry unexplained, can't straightforwardly handle all architectures, and decomposes activations but not the mechanisms themselves
\item SDL latents are dataset-dependent and may not contain the concepts needed for downstream use
\item Lack of solid theoretical foundations for decomposition overall (no agreed-upon definition of ``feature'')
\item Intrinsic interpretability (training models to be decomposable) hasn't yet been competitive
\end{itemize}
}
\end{frame}

\begin{frame}{Open Problems in MechInterp}
Description: interpreting component roles
\begin{itemize}
\item Highly activating examples are correlational
\item Attribution methods are first-order approximations
\item Causal interventions are expensive at scale
\end{itemize}
\source{sharkey2025}

\note{
\begin{itemize}
\item Highly activating examples are correlational and prone to interpretability illusions
\item Attribution methods are theoretically limited (first-order approx) and can be adversarially fooled
\item Feature synthesis struggles with real-world models
\item Causal interventions are expensive at scale; gradient approximations are only approximate
\end{itemize}
}
\end{frame}

\begin{frame}{Open Problems in MechInterp}
Validation: testing our explanations
\begin{itemize}
\item Different methods yield different explanations
\item Performance is seldom demonstrated on real-world tasks
\item Conclusions are often not validated
\end{itemize}
\source{sharkey2025}

\note{
\begin{itemize}
\item Hypotheses are routinely conflated with conclusions
\item Different methods yield different interpretations for the same phenomenon
\item Interpretability methods rarely demonstrate competitive performance on real-world tasks
\item Lack of consensus model organisms and standardized benchmarks
\end{itemize}
}
\end{frame}

\begin{frame}{Open Problems in MechInterp}
Automation: scaling up interpretability
\begin{itemize}
\item Scaling up current methods does not yield more interesting results
\item ``Streetlight interpretability''
\end{itemize}
\source{sharkey2025}

\note{
\begin{itemize}
\item Circuit discovery has low faithfulness, struggles with backup/negative behavior, and is biased toward simple tasks (``streetlight interpretability'')
\item Fully automating current pipelines would still not yield satisfactory explanations
\end{itemize}
}
\end{frame}

\begin{frame}{Open Problems in MechInterp}
Applications: using interpretability in practice
\begin{itemize}
\item Can't reliably distinguish recognizing vs.\ applying a concept
\item Mechanistic verification of behavior not possible
\item Mechanistic understanding might further capabilities by pruning models
\end{itemize}
\source{sharkey2025}

\note{
\begin{itemize}
\item Monitoring/auditing: can't yet reliably distinguish features that recognize deception from those that cause it; mechanistic anomaly detection needs better decomposition
\item Control: activation steering, unlearning, and model editing all limited by inability to isolate individual mechanisms; finetuning for safety is shown to be shallow and easily reversed
\item Predicting behavior: formal verification of AI is far beyond current capabilities; predicting emergent capabilities during training requires understanding dynamic mechanistic changes, not just static snapshots
\item Improving inference/training: mechanistic understanding could enable pruning, distillation, and data selection, but is underexplored
\item Microscope AI: potential for scientific knowledge discovery, but requires rare cross-disciplinary expertise
\item Broader model families: most work is on transformers; unclear how methods/conclusions transfer to SSMs, diffusion models, etc.
\end{itemize}
}
\end{frame}

\begin{frame}{Critiques of Mechanistic Interpretability}
Pick one paper to read and discuss:
\vspace{1em}
\begin{enumerate}
\setlength{\itemsep}{0.5em}
\item \emph{A Pragmatic Vision for Interpretability} (Nanda et al., 2025) \hfill $\sim$15 min
\item \emph{Against Almost Every Theory of Impact of Interpretability} (Ségerie, 2023) \hfill $\sim$20 min
\item \emph{The Misguided Quest for Mechanistic AI Interpretability} (Hendrycks, 2025) \hfill $\sim$20 min
\item \emph{Activation Space Interpretability May Be Doomed} (Chughtai, 2025) \hfill $\sim$15 min
\end{enumerate}

\note{
\begin{itemize}
\item Read one of these papers
\item Come together with people who have read the same paper, discuss the critique: do you agree or disagree?
\item Afterwards, we come back together and discuss the things you found out
\end{itemize}
}
\end{frame}
